\subsection{21.76: closed irreducible nets without a diagonal completion}
\label{cand:21.76}
\problemstatement{21.76}
For every field \(k\) and every \(n\ge3\), set
\[
 R=k[x,y],\quad K=k(x,y),\quad J=yR,\qquad
 \sigma_{12}=\sigma_{21}=kx+J,\quad
 \sigma_{ij}=J\text{ otherwise }(i\ne j).
\]
This gives an irreducible closed elementary net over \(K\) which
cannot be completed by diagonal additive groups.
In particular it supplies every odd characteristic requested.

Recall that an elementary net requires
\(\sigma_{ik}\sigma_{kj}\subseteq\sigma_{ij}\) for pairwise distinct
indices; irreducibility here means every level is nonzero.
All levels contain \(J\), and at least one of \((i,k),(k,j)\)
is outside the two exceptional positions, so every required product
lies in \(RJ=J\). Thus it is an irreducible elementary net.

We first check the rank-one pair. In \(\mathrm{SL}_2(k(x))\), let
\[
 H=\left\langle
 U(a)=\begin{pmatrix}1&ax\\0&1\end{pmatrix},
 V(b)=\begin{pmatrix}1&0\\bx&1\end{pmatrix}:a,b\in k
 \right\rangle.
\]
Its intersection with either elementary root subgroup has parameters
exactly \(kx\).
To see this, combine adjacent factors of the same kind and remove
zero parameters. Any remaining word is alternating with all
parameters nonzero. If its length \(m\ge2\) is even and it starts
with \(U\), choosing all \(m\) off-diagonal terms gives the unique
degree-\(m\) contribution in entry \((1,1)\), nonzero because
\(E_{12}E_{21}\cdots E_{21}=E_{11}\).
If \(m\) is odd, a degree-\((m-1)\) contribution omits one factor.
Omitting the last gives \(E_{11}\), omitting the first gives
\(E_{22}\), and omitting any interior factor gives zero by adjacent
equal matrix units. Thus entry \((1,1)\) has nonzero leading term
of degree \(m-1\). For a word starting with \(V\), interchange the
indices. No such word is an elementary transvection, whose diagonal
entries are both one. Length one and zero give exactly the stated
intersections.
This is a special case of Koibaev's prior closed-pair theorem
\cite{koibaev-closed-pairs}; its needed special case has just been proved.

Let \(E(\sigma)\) be generated by \(I+aE_{ij}\) with
\(a\in\sigma_{ij}\).
All these matrices and their inverses lie in \(\mathrm{SL}_n(R)\).
If \(I+aE_{ij}\in E(\sigma)\), initially with \(a\in K\), its
matrix entry therefore gives \(a\in R\).
Reduction at \(y=0\), now a legitimate ring map on \(R\), sends
\(E(\sigma)\) to \(\operatorname{diag}(H,I_{n-2})\).
At an ordinary position this forces \(a\bmod y=0\); at either
exceptional position the rank-one result gives \(a\bmod y\in kx\).
In each case \(a\in\sigma_{ij}\), proving closedness.
No evaluation homomorphism on the entire field \(K\) was used.

A diagonal completion would force
\(x^2\in\sigma_{11}\) from the two exceptional levels, and then
\(x^3\in\sigma_{11}\sigma_{12}\subseteq\sigma_{12}\).
Reduction modulo \(y\) would give \(x^3\in kx\), impossible for
transcendental \(x\).
Thus no completion exists.

The archive also gives a one-variable characteristic-three example.
In \(\F_9=\F_3[j]\), \(j^2=-1\), the two elementary matrices with
upper parameter \(1\) and lower parameter \(j\) generate a group
of order \(120\), established by an explicit finite coset table.
Its intersections with the order-nine root groups therefore have
parameters \(\F_3\) and \(j\F_3\).
Using exceptional levels
\(\F_3+t\F_9[t]\), \(j\F_3+t\F_9[t]\) and all other levels
\(t\F_9[t]\) repeats the reduction argument over \(\F_9(t)\);
a completion would force \(j\) into the first exceptional level.
The full table is \artifact{research/21.76-finite-table.md}.
This supplementary finite certificate is not needed for the
two-variable proof, which is independent of the carpet candidates
in Problems~19.61--19.62.
See \artifact{research/21.76-proof.md}.
